\chapter{The Two Past Tenses}
\label{ch:perfekt}

German has two ways of talking about the past: the \emph{Perfekt}, built from
\emph{haben} plus a wrapped-up participle, and the \emph{Pr\"{a}teritum}, a
single word with a \emph{-te} on the end. They mean the same thing. What
separates them is where you are and whether you are writing.
\emph{Practice lessons: \texttt{lessons/GE-C15-*.md}.}

\section{\emph{ich habe gesagt} --- the past built from ``have''}
\label{lesson:perfekt}

German's everyday past tense works almost exactly like the English perfect ---
except for two things German does that English doesn't.

\begin{morphologybox}
The recipe: \textbf{haben} (conjugated) $+$ \textbf{past participle}, and the
participle goes \textbf{last}.
\end{morphologybox}

\begin{center}
\begin{tabular}{@{}ll@{}}
\toprule
Infinitive & Participle \\
\midrule
\emph{sag\textbf{en}} & \textbf{ge}\emph{sag}\textbf{t} \\
\emph{mach\textbf{en}} & \textbf{ge}\emph{mach}\textbf{t} \\
\emph{lern\textbf{en}} & \textbf{ge}\emph{lern}\textbf{t} \\
\emph{wohn\textbf{en}} & \textbf{ge}\emph{wohn}\textbf{t} \\
\bottomrule
\end{tabular}
\end{center}

Regular (weak) verbs take \textbf{ge-} in front and \textbf{-t} behind. The
participle is \textbf{wrapped}, not merely suffixed --- what grammarians call a
\emph{circumfix}.

\begin{center}
\begin{tabular}{@{}ll@{}}
\toprule
German & English \\
\midrule
\emph{ich \textbf{habe} gesagt} & I said \\
\emph{du \textbf{hast} gesagt} & you said \\
\emph{wir \textbf{haben} gesagt} & we said \\
\bottomrule
\end{tabular}
\end{center}

\begin{grammarlens}
\textbf{1. The participle goes to the end.} Not next to the verb --- the
\emph{end of the clause}:\\[3pt]
\emph{Ich habe gestern Deutsch \textbf{gelernt}.}\\
(word for word: ``I have yesterday German \textbf{learned}.'')\\[3pt]
Everything else piles up in the middle, and the participle waits at the back.
English cannot do this: ``I have yesterday German learned'' is not a
sentence.\\[3pt]
\textbf{2. It means the plain past.} \emph{Ich habe gesagt} translates as
``\textbf{I said},'' not only ``I have said.'' Where English keeps \emph{I said}
and \emph{I have said} apart, spoken German mostly uses the \emph{Perfekt} for
both.
\end{grammarlens}

\subsection*{The \emph{ge-} that English lost}

\begin{cousinweb}
\textbf{ge-} $\leftarrow$ Germanic \emph{*ga-}, a prefix meaning
``\textbf{together, completely}.'' It marked an action as \textbf{finished} ---
perfective --- which is exactly what a past participle is for. English had it
too, as \emph{y-}, and then dropped it. Two fossils survive:
\end{cousinweb}

\begin{center}
\begin{tabular}{@{}ll@{}}
\toprule
Word & Hiding inside \\
\midrule
\textbf{enough} & Old English \emph{ge}\emph{n\=og} \\
\emph{yclept} (archaic, ``named'') & \emph{ge}\emph{cleopod} \\
\bottomrule
\end{tabular}
\end{center}

So English once wrapped its participles the same way, and German simply never
stopped. Say the wrap out loud --- \textbf{ge}-\emph{sag}-\textbf{t} --- and
then the fossil: \emph{en-\textbf{ough}} $\leftarrow$ \emph{gen\=og}.

\section{\emph{ich sagte} --- the simple past}
\label{lesson:praeteritum}

German has a \textbf{second} past tense, the \emph{Pr\"{a}teritum} --- and
choosing between the two is about \textbf{where you are} and \textbf{whether you
are writing}, not about meaning.

\begin{center}
\begin{tabular}{@{}lll@{}}
\toprule
& \emph{Pr\"{a}teritum} & \emph{Perfekt} \\
\midrule
I said & \emph{\textbf{ich sagte}} & \emph{ich habe gesagt} \\
I made & \emph{\textbf{ich machte}} & \emph{ich habe gemacht} \\
I had & \emph{\textbf{ich hatte}} & \emph{ich habe gehabt} \\
\bottomrule
\end{tabular}
\end{center}

Weak verbs add \textbf{-te}. Same meaning as the \emph{Perfekt}; different
register.

\begin{morphologybox}
That \textbf{-te} is the \textbf{dental preterite} --- the Germanic way of making
a past tense by adding a \emph{t}/\emph{d} sound:\\[3pt]
German: \emph{sag} $\to$ \emph{sag\textbf{te}}\\
English: \emph{walk} $\to$ \emph{walk\textbf{ed}}\\
Dutch: \emph{werk} $\to$ \emph{werk\textbf{te}}\\[3pt]
This is a \textbf{Germanic invention}. Latin had nothing like it --- Romance
builds its past from inherited perfect endings instead. So \emph{sagte} and
\emph{walked} are the same machinery, and French \emph{parla} or Spanish
\emph{habl\'{o}} are a completely different one.
\end{morphologybox}

\subsection*{Who says which}

\begin{center}
\begin{tabular}{@{}ll@{}}
\toprule
Tense & Where it lives \\
\midrule
\emph{Perfekt} & everyday speech --- dominant, especially in the \textbf{south} \\
\emph{Pr\"{a}teritum} & rarer in speech; \textbf{standard in narrative writing} \\
\bottomrule
\end{tabular}
\end{center}

The \emph{Pr\"{a}teritum} has largely vanished from conversation in southern
Germany, Austria and Switzerland. Further north it survives better; in books,
it is the normal past tense. A few verbs resist everywhere: \emph{war},
\emph{hatte}, and the modal verbs stay common even in speech.

\subsection*{The same retreat, three times}

\begin{center}
\begin{tabular}{@{}lll@{}}
\toprule
Language & Simple past & Replaced in speech by \\
\midrule
\textbf{German} & \emph{Pr\"{a}teritum} & \emph{Perfekt} \\
French & \emph{pass\'{e} simple} & \emph{pass\'{e} compos\'{e}} \\
Italian & \emph{passato remoto} & \emph{passato prossimo} \\
\bottomrule
\end{tabular}
\end{center}

\begin{culture}
Two Romance languages and one \textbf{Germanic} one, doing the same thing over
the same centuries --- and \textbf{not} by coincidence. Southern Germany, France
and northern Italy are \textbf{neighbours}, and the change spread through that
connected block by \textbf{contact}. Linguists call this an \emph{areal} change:
it travels across language boundaries because the speakers do, which is why a
Germanic language and two Romance ones line up here while Spanish and
Portuguese, out at the \textbf{western edge}, largely kept their simple past.
It also explains the map \emph{inside} German: the \emph{Pr\"{a}teritum} is
weakest in the \textbf{south}, nearest the centre of the change, and strongest
in the north.
\end{culture}

Practise the pair back to back --- \emph{ich \textbf{sagte}} \ldots{} \emph{ich
\textbf{habe gesagt}}, \emph{Er machte das} \ldots{} \emph{Er hat das gemacht}
--- and hear that nothing about the meaning has moved. Next: the verbs that
build their perfect with \emph{sein}, not \emph{haben}.
